\chapter{Installation}

\label{ch:install}

pii installs as a \emph{user} service by default. There is no package and, in that scope,
no root-owned state: the binary, its configuration, its graphs and its stores all live under
your own home, which is also what makes an upgrade a matter of replacing files that the
installer has already moved aside. A system scope exists --- \texttt{SCOPE=system}, with
\texttt{/etc/pii} and \texttt{/var/lib/pii} --- and everything below is written to be read
with either in mind. This chapter is the whole install path --- the two
directories, the script that fills them, the unit that supervises the process, the
credential file, and the commands that prove the result. It also lists the
preconditions the installer checks before it changes anything, because every one of them
was added after an install that half-succeeded.

\section{What you need}

\begin{itemize}
	\item a Linux host with a writable \texttt{\$HOME} and a user session. systemd
	\texttt{--user} is how the daemon is normally supervised, but it is not required: the
	daemon is a foreground process that can be started by hand, and on a host without a
	user bus that is the only way it runs;
	\item a kernel that can enforce the sandbox profile. Landlock is the mechanism; if the
	profile cannot be enforced the daemon \emph{refuses to start} with exit code 4 rather
	than falling back to unconfined;
	\item the \texttt{pii} binary, built for the host's architecture, plus the three
	sibling binaries if you want them;
	\item \texttt{curl}, which the installer uses for one precondition probe;
	\item Python 3.11 or newer with \texttt{tomllib}, which the installer uses to merge and
	verify the configuration file. The script locates it, or falls back to a \texttt{uv}
	provisioned interpreter under \texttt{\_tmp/};
	\item a provider credential for any workflow that calls a model. It is kept in the
	environment of the service, never in a file that is committed;
	\item a \emph{deploy card} for the revision you are about to install. It is a JSON file
        under \texttt{\_tmp/deploy/} named \texttt{card-\textless{}rev\textgreater{}.json},
keyed by the seven- or
	twelve-character short revision of \texttt{HEAD}. The installer refuses without one.
\end{itemize}

The deploy card is the one precondition that surprises people. It is not bureaucracy
around the install script: it is the record that a specific revision was proved red on
its base, green on its branch, gated, and rehearsed for rollback. A card is rendered,
approved or refused with \texttt{pii deploy-card}, and a card missing any of its five
fields is refused with the field named and no approval path printed.

\section{The binary and the state directory}

The binary is self-contained. Its \emph{state} lives in a directory chosen once, by one
function, in this order:

\begin{enumerate}
	\item \texttt{PII\_STATE\_DIR}, when set and non-empty;
	\item otherwise \texttt{\$HOME/.local/state/pii}.
\end{enumerate}

That single resolver is deliberate. Before it, several commands each had their own notion
of where the store lived, and a command could read a store the daemon was not writing.
The measured symptom was not a crash: it was a CLI reporting an empty run list while the
daemon served hundreds of runs from another file. If you set \texttt{PII\_STATE\_DIR} for
a test, set it for the daemon too, or you are looking at two different systems.

The \emph{runtime} directory is a second, separate rule, and mixing the two up is the most
common installation defect:

\begin{enumerate}
	\item \texttt{PII\_RUNTIME\_DIR}, when set;
	\item \texttt{[core].runtime\_dir} from the config file;
	\item \texttt{/run/user/\textless{}uid\textgreater{}/pii}.
\end{enumerate}

The runtime directory holds the control socket and, on a normal host, is a tmpfs. That is
why the durable stores --- the channel cursor and the timer wheel --- default to
\texttt{\textless{}state dir\textgreater{}/durable} and never to the runtime directory: a timer
that disappears at
reboot is worse than no timer at all. It is also why the installer symlinks the graphs
directory into the runtime directory on every start rather than copying it there.

\subsection*{The shape of an installation}

\begin{center}
\resizebox{\linewidth}{!}{
\begin{tikzpicture}[node distance=10mm and 14mm]
  \node[b] (unit) {unit file\\\scriptsize \texttt{\$UNIT\_DIR/pii.service}};
  \node[b, right=of unit] (bin) {binary\\\scriptsize \texttt{\$PREFIX/pii}};
  \node[b, right=of bin] (conf) {config\\\scriptsize \texttt{\$CONFIG\_DIR/pii.toml}};
  \node[b, right=of conf] (env) {secrets\\\scriptsize \texttt{pii.env}, 0600};
  \node[s, below=of bin] (state) {state dir\\\scriptsize stores, logs, graphs};
  \node[s, right=of state] (run) {runtime dir\\\scriptsize socket, graphs link};
  \node[s, below=of state] (dur) {durable/\\\scriptsize channel, timer};
  \draw[f] (unit) -- node[p,above=2pt,midway]{starts} (bin);
  \draw[f] (bin) -- node[p,above=2pt,midway]{reads} (conf);
  \draw[sf] (conf) -- node[p,above=2pt,midway]{same dir} (env);
  \draw[f] (bin.south) -- (state.north);
  \draw[f] (bin.south east) -- (run.north);
  \draw[sf] (state) -- (dur);
\end{tikzpicture}}
\captionof{figure}{What an installation owns, in the user scope. Notice the two dashed edges:
the key file sits beside the config but is read by the service, and the unit is only a file
until an operator enables it.}
\end{center}

\section{The installer, stage by stage}

One command does the whole install, and it is idempotent. It never deletes: anything it
replaces is moved into a dated directory first, so a failed install leaves you able to
compare and to roll back.

\begin{lstlisting}[language=bash]
bash deploy/install-pii.sh                    # SCOPE=user (the default)
SCOPE=system bash deploy/install-pii.sh       # /etc/systemd/system, /etc/pii, /var/lib/pii
CONFIG_ONLY=1 bash deploy/install-pii.sh      # rewrite the config and credentials only
\end{lstlisting}

The stages run in a fixed order, and the preconditions are all checked before the first
file moves. That ordering is not cosmetic: it means an aborted install changes nothing,
which is the only property that makes a script safe to run on a host you care about. The
one exception is \texttt{CONFIG\_ONLY=1}, which skips the build, the card, the binaries and
the unit, and touches only the config file and the credential file.

\begin{center}
\resizebox{\linewidth}{!}{
\begin{tikzpicture}[node distance=8mm and 12mm]
  \node[b] (pre) {preconditions\\\scriptsize card, model host, prefix};
  \node[b, right=of pre] (build) {build\\\scriptsize clean export};
  \node[b, right=of build] (move) {move aside\\\scriptsize dated directory};
  \node[b, right=of move] (install) {install\\\scriptsize binaries, config, unit};
  \node[v, right=of install] (ok) {refused?};
  \node[b, above=of ok] (abort) {abort\\\scriptsize nothing changed};
  \node[b, below=of ok] (print) {operator steps\\\scriptsize nothing enabled};
  \draw[f] (pre) -- node[p,above=2pt,midway]{all met} (build);
  \draw[f] (build) -- (move);
  \draw[f] (move) -- (install);
  \draw[f] (install) -- (ok);
  \draw[bk] (ok) -- node[p,left=2pt,midway]{yes} (abort);
  \draw[f] (ok) -- node[p,left=2pt,midway]{no} (print);
\end{tikzpicture}}
\captionof{figure}{The installer's shape. Every refusal happens either before the build or in
the config-merge step, which restores the replaced file byte for byte. The script ends by
printing what an operator still has to do; it never enables or starts the unit.}
\end{center}

\subsection*{What the preconditions cost when they fail}

Each precondition exists because the failure it prevents was measured.

\begin{itemize}
	\item \textbf{The model host must answer.} The script does
        \texttt{GET \textless{}llama host\textgreater{}/models} with a five-second limit and
aborts if it does not.
	Without this probe an install succeeds and every graph that needs a model fails at
	spawn, minutes later, with the daemon looking healthy.
	\item \textbf{The named graphs must exist.} The generated configuration names
	\texttt{discord-triage} and \texttt{discord-digest}; the script refuses if either file
	is missing from \texttt{graphs/}, because a config that names a graph the daemon cannot
	load is a daemon that starts with a channel and no work.
	\item \textbf{The prefix must be writable.} It is created if absent and refused if it
	cannot be written --- before the build, so a permissions problem costs seconds rather
	than a release compile.
	\item \textbf{The working tree must be a git checkout.} The card is keyed by the short
	revision of \texttt{HEAD}, and there is no install without a card.
\end{itemize}

The build itself is worth one paragraph, because two deployment defects are recorded in
its comments. The script exports \texttt{git archive HEAD} into a temporary directory and
builds there. Measured on three consecutive deploys, building in the working tree shipped
binaries that reported \texttt{\textless{}rev\textgreater{}-dirty} whenever any path was
uncommitted --- including
a worker's leftover --- so the revision the deploy card quotes carried a suffix meaning
\emph{this is not the commit you verified}. The export has no \texttt{.git}, so
\texttt{build.rs} cannot read the revision from it; the script therefore passes
\texttt{PII\_BUILD\_REVISION} explicitly. That was measured too: the first export-built
binary reported \texttt{unknown} for a build correctly made from a known commit, and
\texttt{pii health} classified it as unverified for ever after.

\section{Environment and scope}

Every path the installer writes is an environment variable with a default, so a second
installation on the same host --- a test prefix, a different user --- needs no editing of
the script.

\begin{longtable}{@{}p{0.20\textwidth}p{0.36\textwidth}p{0.40\textwidth}@{}}
	\toprule
	\textbf{Variable} & \textbf{User scope default} & \textbf{System scope default} \\
	\midrule
	\endhead
	\texttt{SCOPE} & \texttt{user} & \texttt{system} \\
	\texttt{PREFIX} & \texttt{\$HOME/.local/bin} & \texttt{\$HOME/.local/bin} (the static
	unit in \texttt{deploy/} instead names \texttt{/usr/local/bin/pii}) \\
	\texttt{CONFIG\_DIR} & \texttt{\$HOME/.config/pii} & \texttt{/etc/pii} \\
	\texttt{STATE\_DIR} & \texttt{\$HOME/.local/state/pii} & \texttt{/var/lib/pii} \\
	\texttt{RUNTIME\_DIR} & \texttt{/run/user/\textless{}uid\textgreater{}/pii} & \texttt{/run/pii} \\
	\texttt{UNIT\_DIR} & \texttt{\$HOME/.config/systemd/user} & \texttt{/etc/systemd/system} \\
	\texttt{BIN\_SRC} & \texttt{\textless{}repo\textgreater{}/target/release} & same \\
	\texttt{LLAMA\_HOST} & \texttt{http://10.10.10.217:8082/v1} & same \\
	\bottomrule
\end{longtable}

Three more variables change what the script does rather than where it writes.
\texttt{SKIP\_BUILD=1} installs the binaries already in \texttt{BIN\_SRC} and is how a
host without a Rust toolchain takes a release;
\texttt{CONFIG\_ONLY=1} rewrites only the configuration and the channel credentials, with
no build, no card, no binaries and no unit, which is the supported way to change a model
host or a token on a running installation; and \texttt{PII\_PYTHON} pins the interpreter
used for the config merge.

\section{The configuration file}

Behaviour lives in one TOML file, and the file is strict: an unknown key is a load
error, so a typo in a sandbox or a runtime path cannot silently become a default. The
installer writes it by \emph{merging} rather than regenerating, and then verifies that
every table and key path of the replaced file is present in the candidate, naming each
miss instead of proceeding. Before that change, a re-install rewrote the file from the
template and carried only \texttt{model} forward --- which is how an operator's sections
disappeared.

The sections the daemon reads are these, and nothing else:

\begin{center}
\begin{tabular}{@{}p{0.22\textwidth}p{0.66\textwidth}@{}}
	\toprule
	\textbf{Section} & \textbf{What it decides} \\
	\midrule
	\texttt{grant\_ceiling} (top level) & the capability names no graph or node may request
	beyond; empty means no ceiling \\
	\texttt{[core]} & model, runtime dir, durable dir, log file and rotation, trajectory
	root and its bounds, the measurement switches \\
	\texttt{[llama]} & the base URL a \texttt{llamacpp-server:} model dials \\
	\texttt{[net]} & egress host patterns and whether private destinations are allowed \\
	\texttt{[memory]}, \texttt{[memory\_graph]} & the memory store's path and graph
	settings \\
	\texttt{[[worker]]} & supervised child processes, with their own in-flight limits \\
	\texttt{[[channel]]} & inbound message sources and the graph each message runs \\
	\texttt{[[mqtt]]} & broker connections, topics and the graph per message \\
	\texttt{[[cron]]} & named schedules, and the graph each one names \\
	\texttt{[hal]}, \texttt{[wasm]}, \texttt{[evolve]} & optional subsystems; absent means
	off \\
	\texttt{[[model\_tier]]} & named model tiers \\
	\bottomrule
\end{tabular}
\end{center}

The \texttt{[core]} keys an administrator actually touches are
\texttt{model}, \texttt{runtime\_dir}, \texttt{durable\_dir}, \texttt{log\_file},
\texttt{log\_rotate\_bytes} (default 10485760), \texttt{log\_keep\_files} (default 5),
\texttt{trajectory\_dir} (absent means trajectories are off),
\texttt{trajectory\_max\_bytes} (default 512 MiB), \texttt{trajectory\_max\_files}
(default 50\,000), and the two measurement switches \texttt{observer} and \texttt{audit},
which default to \texttt{on} and must stay on in production.

\subsection*{The path the CLI reads is not always the path the service reads}

\texttt{pii daemon} takes \texttt{--config}, defaulting to \texttt{\$PII\_CONFIG} and then
\texttt{\$HOME/.pii/pii.toml}. The installer writes \texttt{\$CONFIG\_DIR/pii.toml} ---
\texttt{\$HOME/.config/pii/pii.toml} in the user scope --- and puts
\texttt{PII\_CONFIG} in the unit. The consequence is exact and worth memorising: a shell
that runs \texttt{pii} without \texttt{PII\_CONFIG} set is reading a different file from
the one the daemon loaded, and \texttt{pii health} will show you the daemon's file by
path and hash if you ask. Export \texttt{PII\_CONFIG} in your shell, or accept that the
two are separate.

The same shape applies to the store. \texttt{pii daemon --help} documents the default as
\texttt{\textless{}runtime dir\textgreater{}/store.redb}, while the store commands and the
backup default to the
state directory; the change that would make the daemon agree had not landed at this
revision. Do what the shipped unit does: pass \texttt{--store} explicitly, pointing at
\texttt{\$STATE\_DIR/store.redb}.

\section{Credentials}

Model providers and channels are reached with values taken from the environment of the
process that needs them. Nothing in the configuration file is a secret, and no pii command
prints one: provider errors and events are scrubbed of the key before they are written.

The convention in this deployment is a single environment file next to the config,
owned by the user, mode \texttt{0600}, read by the service unit and never by the
repository:

\begin{lstlisting}[language=bash]
install -m 600 /dev/null ~/.config/pii/pii.env
# then edit it to hold the provider key(s) and the channel ids, and nothing else
chmod 600 ~/.config/pii/pii.env
\end{lstlisting}

The unit reference is written with a leading dash ---
\texttt{EnvironmentFile=-\$CONFIG\_DIR/pii.env} --- which makes the file optional. The
difference matters: without the dash, a missing file is a unit that fails to start, and the
error is about systemd rather than about the credential. With it, the daemon starts and
reports what it does not have --- \texttt{channels: []} --- which is a state to report, not
a failure to diagnose.

\begin{center}
\resizebox{\linewidth}{!}{
\begin{tikzpicture}[node distance=9mm and 13mm]
  \node[b] (ops) {operator\\\scriptsize exports the token};
  \node[b, right=of ops] (envf) {\texttt{pii.env}\\\scriptsize 0600, outside the repo};
  \node[b, right=of envf] (unit) {unit\\\scriptsize \texttt{EnvironmentFile=-}};
  \node[m, right=of unit] (proc) {daemon process\\\scriptsize environment};
  \node[b, below=of proc] (provider) {provider\\\scriptsize one egress host};
  \node[b, below=of envf] (conf) {\texttt{pii.toml}\\\scriptsize never a secret};
  \draw[f] (ops) -- (envf);
  \draw[f] (envf) -- (unit);
  \draw[f] (unit) -- (proc);
  \draw[f] (proc) -- (provider);
  \draw[bk] (conf) -- node[p,left=2pt,midway]{must not} (envf);
\end{tikzpicture}}
\captionof{figure}{The secret path. The value reaches the daemon's environment and nothing else;
the config file names the shape of the connection, never the credential. A key pasted into a
chat or a commit is compromised, and rotation is the only remedy.}
\end{center}

Named variables, so a grep for one of them in the tree is a finding: \texttt{OPENAI\_API\_KEY}
and \texttt{OPENAI\_BASE\_URL}; \texttt{ANTHROPIC\_API\_KEY} and
\texttt{ANTHROPIC\_BASE\_URL}; \texttt{GEMINI\_API\_KEY} and \texttt{GEMINI\_BASE\_URL};
\texttt{PII\_MODEL} for a default model; \texttt{PII\_DISCORD\_TOKEN} and
\texttt{PII\_DISCORD\_CHANNEL\_ID} for the one forum channel the installer configures.
Egress is granted per host: the daemon grants \texttt{http:\textless{}host\textgreater{}} for
the host of the
configured base URL only, and any other host is denied and reported.

\section{Supervising the daemon}

The daemon is a foreground process: \texttt{--foreground} is accepted for supervisors and
changes nothing, because the daemon never forks. The installer writes this unit and does
not enable it:

\begin{lstlisting}
[Unit]
Description=Pii daemon
After=network-online.target

[Service]
Type=simple
$USERLINE
Environment=PII_RUNTIME_DIR=$RUNTIME_DIR
Environment=PII_CONFIG=$CONF
EnvironmentFile=-$CONFIG_DIR/pii.env
WorkingDirectory=$STATE_DIR
ExecStartPre=/usr/bin/env mkdir -p $RUNTIME_DIR
ExecStartPre=/usr/bin/env ln -sfn $STATE_DIR/graphs $RUNTIME_DIR/graphs
ExecStart=$PREFIX/pii daemon --foreground --config $CONF --store $STATE_DIR/store.redb \
  --workspace $STATE_DIR --log-file $STATE_DIR/pii.log \
  --log-rotate-bytes 10485760 --log-keep-files 5
Restart=always
RestartSec=2
TimeoutStopSec=10
NoNewPrivileges=yes

[Install]
WantedBy=$WANTED
\end{lstlisting}

The variables are the installer's own: \texttt{\$USERLINE} is a \texttt{User=} line and is
emitted only in the system scope, and \texttt{\$WANTED} is \texttt{default.target} for a
user unit and \texttt{multi-user.target} for a system one.

Two lines in that unit are load-bearing beyond the obvious. The symlink of the graphs
directory runs on every start because the runtime directory is a tmpfs: the graphs live in
the state directory and are placed where the daemon looks for them at each boot.
\texttt{TimeoutStopSec=10} matches the daemon's own shutdown: it waits at most five seconds
for in-flight blocking work, then exits anyway, so a stop that needs more than ten seconds
is a stuck process and systemd will kill it.

Then enable it, and make it survive logout if the host is meant to run unattended:
\texttt{systemctl --user daemon-reload}, then \texttt{systemctl --user enable --now pii},
then \texttt{loginctl enable-linger} for the account. Nothing in the installer runs those
three commands, which is deliberate --- an install that starts a service is an install you
cannot inspect before it acts.

The repository also ships a static unit for the system scope,
\texttt{deploy/pii.service}, which uses \texttt{DynamicUser=yes} with
\texttt{StateDirectory}, \texttt{RuntimeDirectory} and \texttt{LogsDirectory} so that
\texttt{/var/lib/pii}, \texttt{/run/pii} and \texttt{/var/log/pii} are created, owned and
kept across restarts. Use one or the other, not both: two units for the same daemon is the
same defect as two daemons for the same store.

\section{The graphs an installation serves}

Graphs are installed as files under \texttt{\$STATE\_DIR/graphs}, and the daemon resolves
a name through three directories in order: \texttt{\textless{}runtime dir\textgreater{}/graphs},
then the
\texttt{graphs} directory beside the store, then
\texttt{\$HOME/.local/state/pii/graphs}. A directory listed twice is searched once. That
order is why an installed copy can shadow a repository copy, and why
\texttt{pii graph list} prints a drift column: \texttt{same}, \texttt{DRIFT},
\texttt{installed-only} or \texttt{repo-only} compares each installed file against its
repository twin by the hash of its bytes.

\begin{lstlisting}[language=bash]
pii graph list                  # name, hash, node count, trigger, armed runs, drift
pii graph list --all            # include graphs that are only in the repository
pii graph edit queue --file graphs/queue.toml    # lint, move the old one aside, write atomically
\end{lstlisting}

\texttt{pii graph edit} never deletes a graph: it lints the candidate first, moves the
previous file aside, and writes atomically. A graph that fails lint is not installed at
all, so a bad edit cannot take out a working graph.

\section{Verifying the installation}

Four commands, in increasing depth. Each has an expected reading, and the exit codes are
part of the answer.

\begin{lstlisting}[language=bash]
pii --version                 # the binary you are actually running
pii health                    # 0 = this build, 2 = no daemon, 3 = drift or unverified
pii health --json             # the full sys.status payload, for a script
pii graph list                # does the daemon see the graphs on disk?
\end{lstlisting}

\begin{center}
\resizebox{\linewidth}{!}{
\begin{tikzpicture}[node distance=12mm and 12mm]
  \node[b] (v) {\texttt{pii --version}};
  \node[b, right=of v] (h) {\texttt{pii health}};
  \node[v, right=of h] (code) {exit code?};
  \node[b, right=of code, yshift=16mm] (zero) {0\\\scriptsize same build, same tree};
  \node[b, right=of code] (two) {2\\\scriptsize nothing answered};
  \node[b, right=of code, yshift=-16mm] (three) {3\\\scriptsize drift or unverified};
  \node[b, right=of three] (unit) {\texttt{systemctl --user status pii}};
  \draw[f] (v) -- (h);
  \draw[f] (h) -- (code);
  \draw[f] (code) -- node[p,above=2pt,midway]{0} (zero);
  \draw[f] (code) -- node[p,above=2pt,midway]{2} (two);
  \draw[bk] (code) -- node[p,below=2pt,midway]{3} (three);
  \draw[sf] (two) -- (unit);
\end{tikzpicture}}
\captionof{figure}{The verification ladder, and the three readings of \texttt{pii health}. Exit
3 immediately after an install is correct: the client is the new binary and the daemon is still
the old process. It becomes 0 only after the restart.}
\end{center}

A health answer whose \texttt{stores} list is empty is not a healthy daemon, whatever its
exit code: it is a daemon pointed at an empty directory. Check
\texttt{PII\_STATE\_DIR} and the unit's \texttt{--store} argument before believing anything
else. The same reading applies to \texttt{format\_version}: a store at a version the binary
does not know is refused at open, with the path and both versions in the message, and
nothing is written to the file.

\section{Rolling an installation back}

\texttt{deploy/rollback-pii.sh} undoes the last install without deleting anything. It
stops the unit if it is running, finds the newest
\texttt{\$PREFIX.old-\textless{}epoch\textgreater{}} directory,
and moves its contents back --- binaries to the prefix, \texttt{pii.toml} to the config
directory, the unit to the unit directory, and any \texttt{graph.*} file to the state
directory's \texttt{graphs}. Whatever it replaces is moved to
\texttt{\$PREFIX.rolledback-\textless{}epoch\textgreater{}}, and it prints each file restored.

\begin{lstlisting}[language=bash]
bash deploy/rollback-pii.sh
systemctl --user daemon-reload && systemctl --user restart pii
pii health                                    # back to the older revision, exit 0
\end{lstlisting}

The script is exercised in a scratch prefix by the deployment checks, and the drill is
worth repeating once on a real host before you need it: an install followed by a rollback
followed by \texttt{pii health} returning 0 proves both directions, and costs a minute.

\section{Failure modes of the install itself}

\begin{center}
\begin{tabular}{@{}p{0.34\textwidth}p{0.58\textwidth}@{}}
	\toprule
	\textbf{Symptom} & \textbf{Cause, and what to check} \\
	\midrule
	\texttt{ABORT: no deploy card for revision \ldots} & No
        \texttt{\_tmp/deploy/card-\textless{}rev\textgreater{}.json} for \texttt{HEAD}. Render
and approve a card, or
	install a revision that has one. \\
	\texttt{ABORT: llama.cpp host \ldots\ does not answer} & The model host is down or the
	URL is wrong. Nothing was changed; the probe runs before the build. \\
	\texttt{cargo: command not found} inside the export & A non-interactive shell did not
	carry your \texttt{PATH}. The script now exports it itself; an older copy of the script
	is what fails here. \\
        The installed binary reports \texttt{\textless{}rev\textgreater{}-dirty} & It was built
in the working tree.
	Only a clean-export build is deployable evidence. \\
	The installed binary reports \texttt{unknown} & The export build did not receive
	\texttt{PII\_BUILD\_REVISION}. Rebuild with the script, not by hand. \\
	\texttt{pii health} exits 2 after install & The unit was written but not enabled, or it
	failed to start. \texttt{systemctl --user status pii}, then the log file. \\
	The daemon starts but every graph fails at spawn & The model id is one the provider does
	not serve, or the key is not in the unit's environment. \\
	\bottomrule
\end{tabular}
\end{center}

The last row is the one that looks like a pii bug and usually is not. A model id that the
provider has retired fails at spawn, in a second, with no partial output; the spec graph
takes its \texttt{err} edge to a \texttt{refuse} node, and the queue records the row as
refused rather than failed. Read the id in the graph, then the model the daemon resolved,
then the provider's own list.
